\documentclass[10pt,a4paper]{amsart}
\usepackage[margin=19mm]{geometry}
\usepackage{amsmath,amssymb,mathtools}
\usepackage[scr=rsfs,cal=euler]{mathalfa}
\usepackage{xfrac}
\usepackage[unicode,hidelinks,bookmarks=false]{hyperref}
\newcommand{\ie}{i.e.\ }
\newcommand{\cf}{cf.\ }
\newcommand{\resp}{respectively\ }
\newcommand{\Subsection}{\subsection}
\providecommand{\nobreakdash}{\nobreak\hbox{-}}
\setlength{\emergencystretch}{2em}
\multlinegap0pt
\newtheorem{thm}{Theorem}[section]
\newtheorem{prop}[thm]{Proposition}
\newtheorem{lemm}[thm]{Lemma}
\newtheorem{coro}[thm]{Corollary}
\newtheorem*{rema}{Remark}
\newcommand{\Z}{\mathbb{Z}}
\newcommand{\Aut}{\operatorname{Aut}}
\newcommand{\Hol}{\operatorname{Hol}}
\newcommand{\ci}{\circ}
\begin{document}
\title[Classification of skew left braces with additive group isomo]{Classification of skew left braces with additive group isomorphic to the infinite dihedral group}
\author{Akihide Hanaki, Yuto Sakata, Hiroki Yoshino}
\date{}
\maketitle
\noindent\emph{Source and licence.} Classification of skew left braces with additive group isomorphic to the infinite dihedral group. Version arXiv:2606.29641v2 (CC BY 4.0). This material is distributed under the Creative Commons Attribution 4.0 International licence, http://creativecommons.org/licenses/by/4.0/, as recorded in the arXiv OAI licence metadata for this exact version and shown on the abstract page https://arxiv.org/abs/2606.29641v2. Full rights notice: licenses/arxiv-2606.29641-CC-BY-4.0.txt.\par\smallskip
\noindent\textit{Editorial scope.} Counted unit: the original \texttt{thm} environment \texttt{cls\_gamma} at source lines 669--674 together with its complete proof at lines 676--731; the delivered document keeps source lines 628--732 so that every referenced label and every citation resolves inside it. Native editable source is the paper's own concatenated TeX; the AMS wrapper is editorial formatting only. No publisher class, upstream script or unrelated artwork is redistributed.
\section{Equivalence of $\gamma$-maps}
We determine the equivalences of $\gamma$-maps.

\begin{lemm}\label{lemCindy}
  $\gamma$-maps are characterized by the values at
  $(1,0)$, $(0,1)$, and $(2,0)$.
\end{lemm}

\begin{proof}
  We make a list for the values of $\gamma_{(i)}$. 
  $$
    \begin{array}{c||c|c|c}
      & \gamma_{(i)}(1,0) & \gamma_{(i)}(0,1) &\gamma_{(i)}(2,0)\\
      \hline
      \hline
      (1) & [0,0] & [0,0] & [0,0] \\
      \hline
      (2,k) & [-2,0] & [k,0] & [-4,0]\\
      \hline
      (3) & [-2,0] & [0,1] & [-4,0]\\
      \hline
      (4,k) & [0,0] & [k,1] & [0,0]\\
      \hline
      (5,k) & [2k,1] & [0,0] & [0,0]\\
      \hline
      (6,k) & [2k,1] & [-2k-2,0] & [-4,0]\\
      \hline
      (7,k) & [2k,1] & [2k,1] & [0,0]\\
      \hline
      (8,k) & [2k,1] & [0,1] & [-4,0]\\
      \hline
    \end{array}
  $$
  This shows the lemma.
\end{proof}

\begin{rema}
  Cindy Tsang pointed out that Lemma \ref{lemCindy}
  holds without the classification of $\gamma$-maps.
\end{rema}


\begin{thm}\label{cls_gamma}
  Representatives of $\gamma$-maps on $D_\infty$ are
  $$\gamma_{(1)}, \gamma_{(2,0)}, \gamma_{(2,1)}, \gamma_{(3)},
  \gamma_{(4,0)}, \gamma_{(4,1)}, \gamma_{(5,0)}, \gamma_{(5,1)},
  \gamma_{(6,0)}, \gamma_{(6,1)}. $$
\end{thm}

\begin{proof}
  We will show that
  \begin{itemize}
    \item $\{\gamma_{(1)}\}$ and $\{\gamma_{(3)}\}$
    are equivalence classes,
    \item for $i, j\in\{2,4,5,6\}$, $i\ne j$, $\gamma_{(i,k)}$
    and $\gamma_{(j,k')}$ are not equivalent,
    \item for $i\in \{2,4,5,6, 7, 8\}$ and $k, k'\in \Z$,
    $\gamma_{(i,k)}$ and $\gamma_{(i,k')}$ are equivalent
    if and only if $k\equiv k'\pmod{2}$, 
    \item $\gamma_{(5,0)}$ and $\gamma_{(7,1)}$ are equivalent, 
    \item $\gamma_{(5,1)}$ and $\gamma_{(7,0)}$ are equivalent,
    \item $\gamma_{(6,0)}$ and $\gamma_{(8,1)}$ are equivalent, and
    \item $\gamma_{(6,1)}$ and $\gamma_{(8,0)}$ are equivalent.
  \end{itemize}
  These prove the proposition.
  To show them, we will give a table of
  $\gamma_{(i)}^{[b,t]}(a,s)$ for $(a,s)=(1,0), (0,1), (2,0)$.

  {\small
  $$\begin{array}{c||c|c|c||c}
    & \gamma_{(i)}^{[b,t]}(1,0)
    & \gamma_{(i)}^{[b,t]}(0,1)
    & \gamma_{(i)}^{[b,t]}(2,0) & \gamma_{(j)}\\
    \hline
    \hline
    (1) & [0,0] & [0,0] & [0,0] & \gamma_{(1)}\\
    \hline
    (2,k) & [-2,0] & [(-1)^t(k+2b),0] & [-4,0]& \gamma_{(2,(-1)^t(k+2b))}\\
    \hline
    (3) & [-2,0] & [0,1] & [-4,0] & \gamma_{(3)}\\
    \hline
    (4,k) & [0,0] & [(-1)^t(k-2b),1] & [0,0] & \gamma_{(4,(-1)^t(k-2b))}\\
    \hline
    (5,k), \text{$b$:even} & [(-1)^t2(k-b),1] & [0,0] & [0,0] & \gamma_{(5,(-1)^t(k-b))}\\
    (5,k), \text{$b$:odd} & [(-1)^t2(k-b),1]  & [(-1)^t2(k-b),1] & [0,0]
                                & \gamma_{(7,(-1)^t(k-b))}\\
    \hline
    (6,k), \text{$b$:even, $t=0$} & [2(k-b),1] & [-2(k-b)-2,0] & [-4,0] &\gamma_{(6,k-b)}\\
    (6,k), \text{$b$:even, $t=1$} & [-2(k-b+2),1] & [2(k-b+2)-2,0] & [-4,0]&\gamma_{(6,-k+b-2)}\\
    (6,k), \text{$b$:odd, $t=0$} & [2(k-b),1] & [0,1] & [-4,0]&\gamma_{(8,k-b)}\\
    (6,k), \text{$b$:odd, $t=1$} & [-2(k-b+2),1] & [0,1] & [-4,0]&\gamma_{(8,-k+b-2)}\\
    \hline
    (7,k), \text{$b$:even}  & [(-1)^t2(k-b),1] & [(-1)^t2(k-b),1] & [0,0]
                                &\gamma_{(7,(-1)^t(k-b))}\\
    (7,k), \text{$b$:odd}  & [(-1)^t2(k-b),1] & [0,0] & [0,0]
                                &\gamma_{(5,(-1)^t(k-b))}\\
    \hline
    (8,k), \text{$b$:even, $t=0$}& [2(k-b),1] & [0,1] & [-4,0]&\gamma_{(8,k-b)}\\
    (8,k), \text{$b$:even, $t=1$}& [-2(k-b+2),1] & [0,1] & [-4,0]&\gamma_{(8,-k+b-2)}\\
    (8,k), \text{$b$:odd, $t=0$}& [2(k-b),1] & [-2(k-b)-2,0] & [-4,0]&\gamma_{(6,k-b)}\\
    (8,k), \text{$b$:odd, $t=1$}& [-2(k-b+2),1] & [2(k-b+2)-2,0] & [-4,0]&\gamma_{(6,-k+b-2)}\\
    \hline
  \end{array}$$}
  By the table, we can check all required conditions.
\end{proof}


\end{document}
